\section{Wprowadzenie}

\subsection{Sformułowanie problemu}

Długość pobytu w szpitalu jest kluczowym wskaźnikiem w zarządzaniu służbą zdrowia. Dokładne przewidywanie czasu hospitalizacji pozwala na optymalizację wykorzystania łóżek szpitalnych, lepsze planowanie personelu oraz skuteczniejsze zarządzanie kosztami.

Problem przewidywania długości pobytu stanowi zadanie \textbf{regresji}, w którym na podstawie czynników ryzyka zdrowotnego pacjenta należy oszacować liczbę dni hospitalizacji.

\subsection{Cel projektu}

Celem projektu jest opracowanie i ocena modeli sieci neuronowych przewidujących długość pobytu w szpitalu na podstawie danych Healthcare Risk Factors Dataset. Projekt obejmuje przygotowanie danych, implementację trzech architektur (SimpleMLP, TabNet, TabTransformer), systematyczne przeszukiwanie hiperparametrów oraz wybór najlepszego modelu.

\subsection{Wykorzystane metody}

Zastosowano trzy architektury sieci neuronowych dostosowane do danych tabelarycznych:

\begin{itemize}
    \item \textbf{SimpleMLP} -- klasyczna architektura fully-connected (baseline)
    \item \textbf{TabNet} -- architektura z mechanizmem attention i selekcją cech
    \item \textbf{TabTransformer} -- architektura oparta na mechanizmie Transformer
\end{itemize}
